% --- Archon physics formalization source begin ---
% archon:physics
% archon:covers PhyXMiniProblems/problem_phyx_mini_0687.lean
% archon:source-report reports/phyx_mini/problem_phyx_mini_0687.source.json

\chapter{Physics problem phyx\_mini\_0687}
\label{ch:PhyXMiniProblems_problem_phyx_mini_0687}

\paragraph{Problem source.}
\#\# Physical scenario

The 20-g centrifuge at NASA's Ames Research Center in Mountain View, California, is a horizontal, cylindrical tube \textbackslash{}( 58.0 \textbackslash{}text\{ ft\} \textbackslash{}) long and is represented in figure. Assume an astronaut in training sits in a seat at one end, facing the axis of rotation \textbackslash{}( 29.0 \textbackslash{}text\{ ft\} \textbackslash{}) away.

\#\# Question

Determine the rotation rate, in revolutions per second, required to give the astronaut a centripetal acceleration of \textbackslash{}( 20.0g \textbackslash{}).

\#\# Answer choices

- A: A: \textbackslash{}( 64.3 \textbackslash{}text\{ rev/min\} \textbackslash{})

- B: B: \textbackslash{}( 22.5 \textbackslash{}text\{ rev/min\} \textbackslash{})

- C: C: \textbackslash{}( 45.0 \textbackslash{}text\{ rev/min\} \textbackslash{})

- D: D: \textbackslash{}( 36.0 \textbackslash{}text\{ rev/min\} \textbackslash{})

\#\# Figure caption (auxiliary; use the image as primary evidence)

The image depicts a large, horizontal structural beam with triangular truss designs along its length. The beam is supported at the center by a vertical shaft. There are two red clamps beneath the beam. Near the right end, a person is seated in a chair inside the structure. The width of the beam from the center to the right end is labeled as "29 ft."

\#\# Dataset metadata

Rotational Motion; Physical Model Grounding Reasoning, Multi-Formula Reasoning

\paragraph{Recorded answer/context.}
C: C: \textbackslash{}( 45.0 \textbackslash{}text\{ rev/min\} \textbackslash{})

\paragraph{Figure/image path.}
/root/proposal\_for\_physic/science-mango/phyx\_mini\_run/phyx\_data/test\_image/687.png

\paragraph{Formalization target.}
create a compiling Lean file with sorry bodies at `PhyXMiniProblems/problem\_phyx\_mini\_0687.lean`. The Lean declarations must preserve the physical quantities, dimensions or dimensional roles, figure labels, governing-law hypotheses, and final relation expressed by this problem.
Use Mathlib/Physlib names found through LeanExplore where available. If a physics API is missing, introduce faithful local abstractions rather than scalar placeholder aliases.

\begin{theorem}[Physics formalization target]
\label{thm:physics:phyx_mini_0687:target}
\lean{PhyXMiniProblems.ProblemPhyXMini0687.required_rotation_rate_is_choice_C}
\uses{def:physics:phyx-mini-0687:phyxminiproblems-problemphyxmini0687-rotationfrequencyinrevolutionspersecond, def:physics:phyx-mini-0687:phyxminiproblems-problemphyxmini0687-rotationfrequencyinrevolutionsperminute, def:physics:phyx-mini-0687:phyxminiproblems-problemphyxmini0687-amestwentygcentrifugesetup, def:physics:phyx-mini-0687:phyxminiproblems-problemphyxmini0687-matchesamescentrifugescenario, def:physics:phyx-mini-0687:phyxminiproblems-problemphyxmini0687-matchescentrifugeproblemreadouts, def:physics:phyx-mini-0687:phyxminiproblems-problemphyxmini0687-matchesprimarycentrifugefigure, def:physics:phyx-mini-0687:phyxminiproblems-problemphyxmini0687-usesroundedstandardgravity, def:physics:phyx-mini-0687:phyxminiproblems-problemphyxmini0687-hasphysicalcentrifugeparameters, def:physics:phyx-mini-0687:phyxminiproblems-problemphyxmini0687-satisfiesstatedglevel, def:physics:phyx-mini-0687:phyxminiproblems-problemphyxmini0687-satisfiesuniformcircularmotion, def:physics:phyx-mini-0687:phyxminiproblems-problemphyxmini0687-answerchoice, def:physics:phyx-mini-0687:phyxminiproblems-problemphyxmini0687-answerchoice-raterevolutionsperminute}
The assigned autoformalize agent should translate this physics problem into Lean declarations in the covered file, with theorem and lemma proof bodies written as `by sorry`.
\end{theorem}
\begin{proof}
This is an autoformalization task, not a proof task. Produce statements that can later be proved by the physics prover without weakening the physical model.
\end{proof}

% --- Archon declaration topology begin ---
\section{Declaration topology}
\begin{definition}[acceleration Dimension]
  \label{def:physics:phyx-mini-0687:phyxminiproblems-problemphyxmini0687-accelerationdimension}
  \lean{PhyXMiniProblems.ProblemPhyXMini0687.accelerationDimension}
  The physical dimension of acceleration, length / time\textasciicircum{}2
\end{definition}
\begin{proof}
  This is the declared model definition of acceleration Dimension.
\end{proof}
\begin{definition}[inverse Time Dimension]
  \label{def:physics:phyx-mini-0687:phyxminiproblems-problemphyxmini0687-inversetimedimension}
  \lean{PhyXMiniProblems.ProblemPhyXMini0687.inverseTimeDimension}
  The physical dimension of both angular speed and rotation frequency
\end{definition}
\begin{proof}
  This is the declared model definition of inverse Time Dimension.
\end{proof}
\begin{definition}[Length Quantity]
  \label{def:physics:phyx-mini-0687:phyxminiproblems-problemphyxmini0687-lengthquantity}
  \lean{PhyXMiniProblems.ProblemPhyXMini0687.LengthQuantity}
  A nonnegative physical length, independent of the selected unit
\end{definition}
\begin{proof}
  This is the declared model definition of Length Quantity.
\end{proof}
\begin{definition}[Acceleration Quantity]
  \label{def:physics:phyx-mini-0687:phyxminiproblems-problemphyxmini0687-accelerationquantity}
  \lean{PhyXMiniProblems.ProblemPhyXMini0687.AccelerationQuantity}
  \uses{def:physics:phyx-mini-0687:phyxminiproblems-problemphyxmini0687-accelerationdimension}
  A nonnegative physical acceleration magnitude
\end{definition}
\begin{proof}
  This is the declared model definition of Acceleration Quantity.
\end{proof}
\begin{definition}[Angular Speed Quantity]
  \label{def:physics:phyx-mini-0687:phyxminiproblems-problemphyxmini0687-angularspeedquantity}
  \lean{PhyXMiniProblems.ProblemPhyXMini0687.AngularSpeedQuantity}
  \uses{def:physics:phyx-mini-0687:phyxminiproblems-problemphyxmini0687-inversetimedimension}
  A nonnegative angular-speed magnitude, with radians dimensionless
\end{definition}
\begin{proof}
  This is the declared model definition of Angular Speed Quantity.
\end{proof}
\begin{definition}[Rotation Frequency Quantity]
  \label{def:physics:phyx-mini-0687:phyxminiproblems-problemphyxmini0687-rotationfrequencyquantity}
  \lean{PhyXMiniProblems.ProblemPhyXMini0687.RotationFrequencyQuantity}
  \uses{def:physics:phyx-mini-0687:phyxminiproblems-problemphyxmini0687-inversetimedimension}
  A nonnegative rotation frequency, with revolutions dimensionless
\end{definition}
\begin{proof}
  This is the declared model definition of Rotation Frequency Quantity.
\end{proof}
\begin{definition}[length Readout]
  \label{def:physics:phyx-mini-0687:phyxminiproblems-problemphyxmini0687-lengthreadout}
  \lean{PhyXMiniProblems.ProblemPhyXMini0687.lengthReadout}
  \uses{def:physics:phyx-mini-0687:phyxminiproblems-problemphyxmini0687-lengthquantity}
  Read a physical length in a selected length unit
\end{definition}
\begin{proof}
  This is the declared model definition of length Readout.
\end{proof}
\begin{definition}[acceleration Readout]
  \label{def:physics:phyx-mini-0687:phyxminiproblems-problemphyxmini0687-accelerationreadout}
  \lean{PhyXMiniProblems.ProblemPhyXMini0687.accelerationReadout}
  \uses{def:physics:phyx-mini-0687:phyxminiproblems-problemphyxmini0687-accelerationquantity}
  Read an acceleration in coherent units induced by length and time units
\end{definition}
\begin{proof}
  This is the declared model definition of acceleration Readout.
\end{proof}
\begin{definition}[angular Speed Readout]
  \label{def:physics:phyx-mini-0687:phyxminiproblems-problemphyxmini0687-angularspeedreadout}
  \lean{PhyXMiniProblems.ProblemPhyXMini0687.angularSpeedReadout}
  \uses{def:physics:phyx-mini-0687:phyxminiproblems-problemphyxmini0687-angularspeedquantity}
  Read an angular speed in radians per selected time unit
\end{definition}
\begin{proof}
  This is the declared model definition of angular Speed Readout.
\end{proof}
\begin{definition}[rotation Frequency Readout]
  \label{def:physics:phyx-mini-0687:phyxminiproblems-problemphyxmini0687-rotationfrequencyreadout}
  \lean{PhyXMiniProblems.ProblemPhyXMini0687.rotationFrequencyReadout}
  \uses{def:physics:phyx-mini-0687:phyxminiproblems-problemphyxmini0687-rotationfrequencyquantity}
  Read a rotation frequency in revolutions per selected time unit
\end{definition}
\begin{proof}
  This is the declared model definition of rotation Frequency Readout.
\end{proof}
\begin{definition}[length In Feet]
  \label{def:physics:phyx-mini-0687:phyxminiproblems-problemphyxmini0687-lengthinfeet}
  \lean{PhyXMiniProblems.ProblemPhyXMini0687.lengthInFeet}
  \uses{def:physics:phyx-mini-0687:phyxminiproblems-problemphyxmini0687-lengthquantity, def:physics:phyx-mini-0687:phyxminiproblems-problemphyxmini0687-lengthreadout}
  Foot readout used for the two stated apparatus lengths
\end{definition}
\begin{proof}
  This is the declared model definition of length In Feet.
\end{proof}
\begin{definition}[length In Meters]
  \label{def:physics:phyx-mini-0687:phyxminiproblems-problemphyxmini0687-lengthinmeters}
  \lean{PhyXMiniProblems.ProblemPhyXMini0687.lengthInMeters}
  \uses{def:physics:phyx-mini-0687:phyxminiproblems-problemphyxmini0687-lengthquantity, def:physics:phyx-mini-0687:phyxminiproblems-problemphyxmini0687-lengthreadout}
  Coherent SI metre readout of a physical length
\end{definition}
\begin{proof}
  This is the declared model definition of length In Meters.
\end{proof}
\begin{definition}[acceleration In Meters Per Second Squared]
  \label{def:physics:phyx-mini-0687:phyxminiproblems-problemphyxmini0687-accelerationinmeterspersecondsquared}
  \lean{PhyXMiniProblems.ProblemPhyXMini0687.accelerationInMetersPerSecondSquared}
  \uses{def:physics:phyx-mini-0687:phyxminiproblems-problemphyxmini0687-accelerationquantity, def:physics:phyx-mini-0687:phyxminiproblems-problemphyxmini0687-accelerationreadout}
  Coherent SI metre-per-second-squared acceleration readout
\end{definition}
\begin{proof}
  This is the declared model definition of acceleration In Meters Per Second Squared.
\end{proof}
\begin{definition}[angular Speed In Radians Per Second]
  \label{def:physics:phyx-mini-0687:phyxminiproblems-problemphyxmini0687-angularspeedinradianspersecond}
  \lean{PhyXMiniProblems.ProblemPhyXMini0687.angularSpeedInRadiansPerSecond}
  \uses{def:physics:phyx-mini-0687:phyxminiproblems-problemphyxmini0687-angularspeedquantity, def:physics:phyx-mini-0687:phyxminiproblems-problemphyxmini0687-angularspeedreadout}
  Angular speed in radians per second
\end{definition}
\begin{proof}
  This is the declared model definition of angular Speed In Radians Per Second.
\end{proof}
\begin{definition}[rotation Frequency In Revolutions Per Second]
  \label{def:physics:phyx-mini-0687:phyxminiproblems-problemphyxmini0687-rotationfrequencyinrevolutionspersecond}
  \lean{PhyXMiniProblems.ProblemPhyXMini0687.rotationFrequencyInRevolutionsPerSecond}
  \uses{def:physics:phyx-mini-0687:phyxminiproblems-problemphyxmini0687-rotationfrequencyquantity, def:physics:phyx-mini-0687:phyxminiproblems-problemphyxmini0687-rotationfrequencyreadout}
  Rotation frequency in revolutions per second
\end{definition}
\begin{proof}
  This is the declared model definition of rotation Frequency In Revolutions Per Second.
\end{proof}
\begin{definition}[rotation Frequency In Revolutions Per Minute]
  \label{def:physics:phyx-mini-0687:phyxminiproblems-problemphyxmini0687-rotationfrequencyinrevolutionsperminute}
  \lean{PhyXMiniProblems.ProblemPhyXMini0687.rotationFrequencyInRevolutionsPerMinute}
  \uses{def:physics:phyx-mini-0687:phyxminiproblems-problemphyxmini0687-rotationfrequencyquantity, def:physics:phyx-mini-0687:phyxminiproblems-problemphyxmini0687-rotationfrequencyreadout}
  Rotation frequency in revolutions per minute
\end{definition}
\begin{proof}
  This is the declared model definition of rotation Frequency In Revolutions Per Minute.
\end{proof}
\begin{lemma}[length In Meters eq feet mul conversion]
  \label{lem:physics:phyx-mini-0687:phyxminiproblems-problemphyxmini0687-lengthinmeters-eq-feet-mul-conversion}
  \lean{PhyXMiniProblems.ProblemPhyXMini0687.lengthInMeters_eq_feet_mul_conversion}
  \uses{def:physics:phyx-mini-0687:phyxminiproblems-problemphyxmini0687-lengthquantity, def:physics:phyx-mini-0687:phyxminiproblems-problemphyxmini0687-lengthinfeet, def:physics:phyx-mini-0687:phyxminiproblems-problemphyxmini0687-lengthinmeters}
  One foot is exactly 0.3048 m = 381/1250 m
\end{lemma}
\begin{proof}
  The conclusion follows from the governing relations and figure readouts named in the statement of length In Meters eq feet mul conversion.
\end{proof}
\begin{lemma}[rotation Frequency per Minute eq sixty mul per Second]
  \label{lem:physics:phyx-mini-0687:phyxminiproblems-problemphyxmini0687-rotationfrequency-perminute-eq-sixty-mul-persecond}
  \lean{PhyXMiniProblems.ProblemPhyXMini0687.rotationFrequency_perMinute_eq_sixty_mul_perSecond}
  \uses{def:physics:phyx-mini-0687:phyxminiproblems-problemphyxmini0687-rotationfrequencyquantity, def:physics:phyx-mini-0687:phyxminiproblems-problemphyxmini0687-rotationfrequencyinrevolutionspersecond, def:physics:phyx-mini-0687:phyxminiproblems-problemphyxmini0687-rotationfrequencyinrevolutionsperminute}
  A per-minute frequency readout is sixty times its per-second readout
\end{lemma}
\begin{proof}
  The conclusion follows from the governing relations and figure readouts named in the statement of rotation Frequency per Minute eq sixty mul per Second.
\end{proof}
\begin{definition}[Tube Orientation]
  \label{def:physics:phyx-mini-0687:phyxminiproblems-problemphyxmini0687-tubeorientation}
  \lean{PhyXMiniProblems.ProblemPhyXMini0687.TubeOrientation}
  Orientation of the long centrifuge tube
\end{definition}
\begin{proof}
  This is the declared model definition of Tube Orientation.
\end{proof}
\begin{definition}[Rotation Axis Location]
  \label{def:physics:phyx-mini-0687:phyxminiproblems-problemphyxmini0687-rotationaxislocation}
  \lean{PhyXMiniProblems.ProblemPhyXMini0687.RotationAxisLocation}
  Location of the rotation axis relative to the full tube
\end{definition}
\begin{proof}
  This is the declared model definition of Rotation Axis Location.
\end{proof}
\begin{definition}[Astronaut Orientation]
  \label{def:physics:phyx-mini-0687:phyxminiproblems-problemphyxmini0687-astronautorientation}
  \lean{PhyXMiniProblems.ProblemPhyXMini0687.AstronautOrientation}
  Direction in which the seated astronaut faces
\end{definition}
\begin{proof}
  This is the declared model definition of Astronaut Orientation.
\end{proof}
\begin{definition}[Figure Object]
  \label{def:physics:phyx-mini-0687:phyxminiproblems-problemphyxmini0687-figureobject}
  \lean{PhyXMiniProblems.ProblemPhyXMini0687.FigureObject}
  Distinct objects visibly present in the supplied centrifuge image
\end{definition}
\begin{proof}
  This is the declared model definition of Figure Object.
\end{proof}
\begin{definition}[Figure Text Label]
  \label{def:physics:phyx-mini-0687:phyxminiproblems-problemphyxmini0687-figuretextlabel}
  \lean{PhyXMiniProblems.ProblemPhyXMini0687.FigureTextLabel}
  The sole numerical distance label visible in the supplied image
\end{definition}
\begin{proof}
  This is the declared model definition of Figure Text Label.
\end{proof}
\begin{definition}[Ames Centrifuge Figure]
  \label{def:physics:phyx-mini-0687:phyxminiproblems-problemphyxmini0687-amescentrifugefigure}
  \lean{PhyXMiniProblems.ProblemPhyXMini0687.AmesCentrifugeFigure}
  \uses{def:physics:phyx-mini-0687:phyxminiproblems-problemphyxmini0687-figureobject, def:physics:phyx-mini-0687:phyxminiproblems-problemphyxmini0687-figuretextlabel}
  Qualitative evidence transcribed from image 687.png. The dimension line runs from the central shaft to the right end containing the astronaut
\end{definition}
\begin{proof}
  This is the declared model definition of Ames Centrifuge Figure.
\end{proof}
\begin{definition}[Ames Twenty GCentrifuge Setup]
  \label{def:physics:phyx-mini-0687:phyxminiproblems-problemphyxmini0687-amestwentygcentrifugesetup}
  \lean{PhyXMiniProblems.ProblemPhyXMini0687.AmesTwentyGCentrifugeSetup}
  \uses{def:physics:phyx-mini-0687:phyxminiproblems-problemphyxmini0687-lengthquantity, def:physics:phyx-mini-0687:phyxminiproblems-problemphyxmini0687-accelerationquantity, def:physics:phyx-mini-0687:phyxminiproblems-problemphyxmini0687-angularspeedquantity, def:physics:phyx-mini-0687:phyxminiproblems-problemphyxmini0687-rotationfrequencyquantity, def:physics:phyx-mini-0687:phyxminiproblems-problemphyxmini0687-tubeorientation, def:physics:phyx-mini-0687:phyxminiproblems-problemphyxmini0687-rotationaxislocation, def:physics:phyx-mini-0687:phyxminiproblems-problemphyxmini0687-astronautorientation, def:physics:phyx-mini-0687:phyxminiproblems-problemphyxmini0687-amescentrifugefigure}
  The requested rotation frequency and its angular speed are independent physical fields. They are related to the supplied geometry and acceleration only through the governing-law structures below; neither is defined using a numerical answer
\end{definition}
\begin{proof}
  This is the declared model definition of Ames Twenty GCentrifuge Setup.
\end{proof}
\begin{definition}[Matches Ames Centrifuge Scenario]
  \label{def:physics:phyx-mini-0687:phyxminiproblems-problemphyxmini0687-matchesamescentrifugescenario}
  \lean{PhyXMiniProblems.ProblemPhyXMini0687.MatchesAmesCentrifugeScenario}
  \uses{def:physics:phyx-mini-0687:phyxminiproblems-problemphyxmini0687-amestwentygcentrifugesetup}
  Qualitative physical scenario stated in the problem prose
\end{definition}
\begin{proof}
  This is the declared model definition of Matches Ames Centrifuge Scenario.
\end{proof}
\begin{definition}[Matches Centrifuge Problem Readouts]
  \label{def:physics:phyx-mini-0687:phyxminiproblems-problemphyxmini0687-matchescentrifugeproblemreadouts}
  \lean{PhyXMiniProblems.ProblemPhyXMini0687.MatchesCentrifugeProblemReadouts}
  \uses{def:physics:phyx-mini-0687:phyxminiproblems-problemphyxmini0687-lengthinfeet, def:physics:phyx-mini-0687:phyxminiproblems-problemphyxmini0687-amestwentygcentrifugesetup}
  Numerical quantities supplied by the problem. No rotation frequency or answer-choice value occurs in this structure
\end{definition}
\begin{proof}
  This is the declared model definition of Matches Centrifuge Problem Readouts.
\end{proof}
\begin{definition}[Matches Primary Centrifuge Figure]
  \label{def:physics:phyx-mini-0687:phyxminiproblems-problemphyxmini0687-matchesprimarycentrifugefigure}
  \lean{PhyXMiniProblems.ProblemPhyXMini0687.MatchesPrimaryCentrifugeFigure}
  \uses{def:physics:phyx-mini-0687:phyxminiproblems-problemphyxmini0687-lengthreadout, def:physics:phyx-mini-0687:phyxminiproblems-problemphyxmini0687-amestwentygcentrifugesetup}
  Direct evidence from the primary bitmap and its printed dimension line
\end{definition}
\begin{proof}
  This is the declared model definition of Matches Primary Centrifuge Figure.
\end{proof}
\begin{definition}[Uses Rounded Standard Gravity]
  \label{def:physics:phyx-mini-0687:phyxminiproblems-problemphyxmini0687-usesroundedstandardgravity}
  \lean{PhyXMiniProblems.ProblemPhyXMini0687.UsesRoundedStandardGravity}
  \uses{def:physics:phyx-mini-0687:phyxminiproblems-problemphyxmini0687-accelerationinmeterspersecondsquared, def:physics:phyx-mini-0687:phyxminiproblems-problemphyxmini0687-amestwentygcentrifugesetup}
  Rounded standard gravitational acceleration used for the classroom calculation. It calibrates g; it does not constrain the rotation rate
\end{definition}
\begin{proof}
  This is the declared model definition of Uses Rounded Standard Gravity.
\end{proof}
\begin{definition}[Has Physical Centrifuge Parameters]
  \label{def:physics:phyx-mini-0687:phyxminiproblems-problemphyxmini0687-hasphysicalcentrifugeparameters}
  \lean{PhyXMiniProblems.ProblemPhyXMini0687.HasPhysicalCentrifugeParameters}
  \uses{def:physics:phyx-mini-0687:phyxminiproblems-problemphyxmini0687-lengthinmeters, def:physics:phyx-mini-0687:phyxminiproblems-problemphyxmini0687-accelerationinmeterspersecondsquared, def:physics:phyx-mini-0687:phyxminiproblems-problemphyxmini0687-angularspeedinradianspersecond, def:physics:phyx-mini-0687:phyxminiproblems-problemphyxmini0687-rotationfrequencyinrevolutionspersecond, def:physics:phyx-mini-0687:phyxminiproblems-problemphyxmini0687-amestwentygcentrifugesetup}
  Positivity assumptions selecting the nondegenerate physical branch
\end{definition}
\begin{proof}
  This is the declared model definition of Has Physical Centrifuge Parameters.
\end{proof}
\begin{definition}[Satisfies Stated GLevel]
  \label{def:physics:phyx-mini-0687:phyxminiproblems-problemphyxmini0687-satisfiesstatedglevel}
  \lean{PhyXMiniProblems.ProblemPhyXMini0687.SatisfiesStatedGLevel}
  \uses{def:physics:phyx-mini-0687:phyxminiproblems-problemphyxmini0687-accelerationinmeterspersecondsquared, def:physics:phyx-mini-0687:phyxminiproblems-problemphyxmini0687-amestwentygcentrifugesetup}
  The requested centripetal acceleration is the stated dimensionless multiple of local standard gravity. This is the meaning of 20 g, not an assumption about the requested rotation rate
\end{definition}
\begin{proof}
  This is the declared model definition of Satisfies Stated GLevel.
\end{proof}
\begin{definition}[Satisfies Uniform Circular Motion]
  \label{def:physics:phyx-mini-0687:phyxminiproblems-problemphyxmini0687-satisfiesuniformcircularmotion}
  \lean{PhyXMiniProblems.ProblemPhyXMini0687.SatisfiesUniformCircularMotion}
  \uses{def:physics:phyx-mini-0687:phyxminiproblems-problemphyxmini0687-lengthinmeters, def:physics:phyx-mini-0687:phyxminiproblems-problemphyxmini0687-accelerationinmeterspersecondsquared, def:physics:phyx-mini-0687:phyxminiproblems-problemphyxmini0687-angularspeedinradianspersecond, def:physics:phyx-mini-0687:phyxminiproblems-problemphyxmini0687-rotationfrequencyinrevolutionspersecond, def:physics:phyx-mini-0687:phyxminiproblems-problemphyxmini0687-amestwentygcentrifugesetup}
  Uniform circular-motion laws in coherent SI readouts. Radians and revolutions are dimensionless, with one revolution corresponding to 2*pi radians
\end{definition}
\begin{proof}
  This is the declared model definition of Satisfies Uniform Circular Motion.
\end{proof}
\begin{lemma}[rotation Frequency sq eq acceleration div radius and four pi sq]
  \label{lem:physics:phyx-mini-0687:phyxminiproblems-problemphyxmini0687-rotationfrequency-sq-eq-acceleration-div-radius-and-four-pi-sq}
  \lean{PhyXMiniProblems.ProblemPhyXMini0687.rotationFrequency_sq_eq_acceleration_div_radius_and_four_pi_sq}
  \uses{def:physics:phyx-mini-0687:phyxminiproblems-problemphyxmini0687-lengthinmeters, def:physics:phyx-mini-0687:phyxminiproblems-problemphyxmini0687-accelerationinmeterspersecondsquared, def:physics:phyx-mini-0687:phyxminiproblems-problemphyxmini0687-rotationfrequencyinrevolutionspersecond, def:physics:phyx-mini-0687:phyxminiproblems-problemphyxmini0687-amestwentygcentrifugesetup, def:physics:phyx-mini-0687:phyxminiproblems-problemphyxmini0687-hasphysicalcentrifugeparameters, def:physics:phyx-mini-0687:phyxminiproblems-problemphyxmini0687-satisfiesuniformcircularmotion}
  An algebraic consequence of the two circular-motion laws. This leaves the positive square-root selection and all numerical evaluation to later proofs
\end{lemma}
\begin{proof}
  The conclusion follows from the governing relations and figure readouts named in the statement of rotation Frequency sq eq acceleration div radius and four pi sq.
\end{proof}
\begin{definition}[Answer Choice]
  \label{def:physics:phyx-mini-0687:phyxminiproblems-problemphyxmini0687-answerchoice}
  \lean{PhyXMiniProblems.ProblemPhyXMini0687.AnswerChoice}
  Labels of the four rotation-rate choices printed by the source
\end{definition}
\begin{proof}
  This is the declared model definition of Answer Choice.
\end{proof}
\begin{definition}[rate Revolutions Per Minute]
  \label{def:physics:phyx-mini-0687:phyxminiproblems-problemphyxmini0687-answerchoice-raterevolutionsperminute}
  \lean{PhyXMiniProblems.ProblemPhyXMini0687.AnswerChoice.rateRevolutionsPerMinute}
  \uses{def:physics:phyx-mini-0687:phyxminiproblems-problemphyxmini0687-answerchoice}
  Revolution-per-minute value displayed beside each answer label
\end{definition}
\begin{proof}
  This is the declared model definition of rate Revolutions Per Minute.
\end{proof}
% --- Archon declaration topology end ---
% --- Archon physics formalization source end ---
